
\documentclass[11pt,spanish]{article}

% =========================================================
% PLANTILLA BASE PARA INFORMES ACADÉMICOS / TÉCNICOS
% =========================================================

% --- Idioma y codificación ---
\usepackage[utf8]{inputenc}
\usepackage[spanish,es-tabla]{babel}
\usepackage[T1]{fontenc}
\usepackage{textcomp}
\usepackage{newunicodechar}
\newunicodechar{₡}{\textcolonmonetary}

% --- Márgenes / papel ---
\usepackage[
    left=2.5cm,
    right=2.5cm,
    top=2.5cm,
    bottom=2.5cm,
    headsep=0.5cm,
    footskip=0.8cm
]{geometry}

% --- Matemáticas ---
\usepackage{amsmath}

% --- Tablas, enlaces y elementos visuales ---
\usepackage{array}
\usepackage{tabularx}
\usepackage{booktabs}
\usepackage{longtable}
\usepackage{graphicx}
\usepackage{float}
\usepackage{xcolor}
\usepackage{tikz}
\usetikzlibrary{positioning,shapes.geometric,arrows.meta,calc}
\usepackage{enumitem}
\usepackage{ragged2e}
\usepackage[hidelinks]{hyperref}

% --- Encabezado y pie de página ---
\usepackage{fancyhdr}
\setlength{\headheight}{14pt}
\usepackage{lastpage}

% --- Interlineado ---
\linespread{1.2}
\sloppy

% =========================================================
% DATOS DEL DOCUMENTO
% =========================================================
\newcommand{\Institucion}{Instituto Tecnológico de Costa Rica}
\newcommand{\Campus}{Campus Tecnológico Central Cartago}
\newcommand{\DocumentoTitulo}{Proceso de Control de Cambios}
\newcommand{\Curso}{Administración de proyectos}
\newcommand{\Profesor}{Alicia Marcela Salazar Hernández}
\newcommand{\FechaDocumento}{18 de junio de 2026}
\newcommand{\PeriodoAcademico}{I Semestre}
\newcommand{\AnioDocumento}{2026}
\newcommand{\AutoresPortada}{%
    \begin{tabular}{l@{\hspace{1.5cm}}r}
        Mauricio González Prendas & 2024143009 \\
        Isaac Jiménez Blanco & 2024202804 \\
        Tamara Robles Camacho & 2024099342 \\
    \end{tabular}%
}
\newcommand{\DocHeaderLeft}{}
\newcommand{\DocHeaderRight}{Proceso de Control de Cambios}
\newcommand{\DocFooterRight}{Página~\thepage\ de~\pageref{LastPage}}

% Metadatos PDF.
\title{\DocumentoTitulo}
\author{\AutoresPortada}
\date{\FechaDocumento}

% =========================================================
% CONFIGURACIÓN DE TABLAS Y UTILIDADES
% =========================================================
\newcolumntype{Y}{>{\RaggedRight\arraybackslash}X}
\newcolumntype{L}[1]{>{\RaggedRight\arraybackslash}p{#1}}
\renewcommand{\arraystretch}{1.22}
\setlength{\LTpre}{0.5em}
\setlength{\LTpost}{0.5em}

% Imagen segura: si el archivo no existe, muestra un recuadro de marcador.
\newcommand{\SafeImage}[2][0.92\textwidth]{%
    \IfFileExists{#2}{%
        \includegraphics[width=#1]{#2}%
    }{%
        \fbox{%
            \begin{minipage}[c][4cm][c]{#1}
            \centering
            Imagen pendiente:\\
            \texttt{#2}
            \end{minipage}%
        }%
    }%
}

% Figura reutilizable.
\newcommand{\EvidenceFigure}[3]{%
    \begin{figure}[H]
    \centering
    \SafeImage{#1}
    \caption{#2}
    \label{#3}
    \end{figure}%
}

% =========================================================
% ENCABEZADO Y PIE
% =========================================================
\pagestyle{fancy}
\fancyhf{}
\fancyhead[L]{\DocHeaderLeft}
\fancyhead[R]{\DocHeaderRight}
\renewcommand{\headrulewidth}{0.4pt}
\renewcommand{\footrulewidth}{0.4pt}
\fancyfoot[R]{\DocFooterRight}

% Estilo equivalente para páginas horizontales.
\fancypagestyle{widefancy}{%
    \fancyhf{}%
    \fancyhead[L]{\DocHeaderLeft}%
    \fancyhead[R]{\DocHeaderRight}%
    \renewcommand{\headrulewidth}{0.4pt}%
    \renewcommand{\footrulewidth}{0.4pt}%
    \fancyfoot[R]{\DocFooterRight}%
}

% =========================================================
% PÁGINAS HORIZONTALES PARA TABLAS ANCHAS
% =========================================================
\makeatletter
\newcommand{\SyncPageBuilderDimensions}{%
    \global\vsize=\textheight
    \global\@colroom=\textheight
    \global\@colht=\textheight
}
\makeatother

\newenvironment{widepage}{%
    \clearpage
    \hypersetup{pageanchor=false}%
    \paperwidth=297mm
    \paperheight=210mm
    \pdfpagewidth=297mm
    \pdfpageheight=210mm
    \newgeometry{left=2.5cm,right=2.5cm,top=2.5cm,bottom=2.5cm,headsep=0.5cm,footskip=0.8cm}%
    \setlength{\textwidth}{243mm}%
    \setlength{\textheight}{156mm}%
    \setlength{\linewidth}{\textwidth}%
    \setlength{\columnwidth}{\textwidth}%
    \hsize=\textwidth
    \setlength{\headwidth}{\textwidth}%
    \SyncPageBuilderDimensions
    \pagestyle{widefancy}%
}{%
    \clearpage
    \paperwidth=210mm
    \paperheight=297mm
    \pdfpagewidth=210mm
    \pdfpageheight=297mm
    \restoregeometry
    \SyncPageBuilderDimensions
    \hypersetup{pageanchor=true}%
    \pagestyle{fancy}%
}

% =========================================================
% PORTADA
% =========================================================
\newcommand{\MakeCover}{%
\begin{titlepage}
\thispagestyle{empty}
\centering

{\bfseries \huge \Institucion \par}
\vspace{0.25cm}
{\LARGE \Campus \par}

\vfill

{\huge \DocumentoTitulo \par}

\vfill

{\bfseries \LARGE Profesora \par}
{\Large \Profesor \par}

\vfill

{\bfseries \LARGE Estudiantes \par}
{\Large \AutoresPortada \par}

\vfill

{\bfseries\LARGE Fecha \par}
{\Large \FechaDocumento \par}

\vfill

{\LARGE \PeriodoAcademico \par}
{\LARGE \AnioDocumento \par}

\end{titlepage}%
}

% =========================================================
% DOCUMENTO
% =========================================================
\begin{document}
\hypersetup{pageanchor=false}

% =========================
% Portada
% =========================
\MakeCover

% =========================
% Índice
% =========================
\pagenumbering{roman}
\pagestyle{empty}
\addtocontents{toc}{\protect\thispagestyle{empty}}
\tableofcontents
\clearpage
\pagenumbering{arabic}
\hypersetup{pageanchor=true}
\pagestyle{fancy}

% =========================
% Contenido
% =========================

\section{Información general del documento}

\begin{table}[H]
\centering
\begin{tabularx}{\textwidth}{p{5cm} Y}
\toprule
\textbf{Nombre del proyecto} & Plataforma Universitaria Integrada \\
\textbf{Organización} & Universidad Tecnológica La Mejor \\
\textbf{Documento} & \DocumentoTitulo \\
\textbf{Directora del proyecto} & Tamara Robles Camacho \\
\textbf{Fecha} & \FechaDocumento \\
\bottomrule
\end{tabularx}
\end{table}

\section{Introducción}

El presente documento define el Proceso de Control de Cambios del proyecto Plataforma Universitaria Integrada. Este proceso tiene como finalidad establecer la forma en que se recibirán, analizarán, aprobarán, rechazarán, documentarán y comunicarán los cambios que puedan afectar el alcance, el cronograma, los costos, la calidad, los riesgos o los entregables del proyecto.

El control de cambios es necesario porque el proyecto contiene varios documentos de gestión y un prototipo web navegable que debe mantenerse alineado con el alcance aprobado. Además, el equipo trabaja con roles combinados y un presupuesto definido mediante estimación PERT, por lo tanto cualquier modificación no controlada puede generar inconsistencias entre documentos, sobrecarga de trabajo, aumento de costos o reducción de calidad.

El proceso descrito en este documento se adapta a la escala académica del proyecto. Por esta razón, no se plantea un comité corporativo complejo, sino un flujo formal y claro donde la directora del proyecto coordina la revisión, los analistas documentan el impacto, los desarrolladores estiman la viabilidad técnica y el patrocinador aprueba los cambios relevantes.

\section{Propósito del proceso}

El propósito del proceso de control de cambios es evitar modificaciones informales que afecten el avance del proyecto sin análisis previo. Asimismo, busca proteger las líneas base de alcance, cronograma, costo y calidad que fueron definidas durante la planificación.

El proceso persigue los siguientes objetivos:

\begin{itemize}[leftmargin=*, label={\textbullet}]
    \item Definir un flujo formal para recibir y evaluar solicitudes de cambio.
    \item Analizar el impacto de cada cambio en alcance, tiempo, costo, calidad y riesgos.
    \item Determinar si un cambio debe aprobarse, rechazarse o posponerse.
    \item Documentar las decisiones para mantener trazabilidad.
    \item Comunicar los cambios aprobados al equipo y a los interesados relevantes.
    \item Actualizar documentos, cronograma o prototipo cuando el cambio sea aprobado.
\end{itemize}

\section{Principios del control de cambios}

El proceso se basa en reglas simples que permiten mantener orden durante la ejecución del proyecto.

\begin{table}[H]
\centering
\caption{Principios del proceso de control de cambios}
\label{tab:principios-cambios}
{\small
\begin{tabularx}{\textwidth}{p{4cm} Y}
\toprule
\textbf{Principio} & \textbf{Aplicación en el proyecto} \\
\midrule
Formalidad & Todo cambio relevante debe documentarse mediante una solicitud de cambio o registro equivalente. \\
Análisis previo & Ningún cambio que afecte alcance, tiempo, costo, calidad o riesgos debe ejecutarse sin revisión de impacto. \\
Trazabilidad & Las decisiones deben quedar respaldadas en documentos, repositorio, correo, minuta o evidencia del proyecto. \\
Coherencia & Si el cambio modifica un documento base, deben actualizarse los entregables relacionados. \\
Responsabilidad & La directora del proyecto coordina la decisión y los responsables técnicos o documentales apoyan el análisis. \\
Control del alcance & No se deben agregar funciones fuera del alcance sin evaluar su efecto sobre el prototipo y el cronograma. \\
\bottomrule
\end{tabularx}
}
\end{table}

\section{Tipos de cambios considerados}

Para este proyecto se consideran cambios en diferentes áreas de gestión y desarrollo. La clasificación permite determinar quién debe revisarlos y qué documentos podrían verse afectados.

\begin{table}[H]
\centering
\caption{Tipos de cambios del proyecto}
\label{tab:tipos-cambios}
{\small
\begin{tabularx}{\textwidth}{p{3.5cm} Y p{3.2cm}}
\toprule
\textbf{Tipo de cambio} & \textbf{Descripción} & \textbf{Ejemplo} \\
\midrule
Alcance & Agrega, elimina o modifica funcionalidades, módulos, entregables o límites del proyecto. & Incorporar una nueva pantalla al Portal Estudiantil. \\
Cronograma & Cambia fechas, duraciones, dependencias, hitos o ruta crítica. & Mover pruebas de QA por atraso en desarrollo. \\
Costo & Afecta el presupuesto, rubros de adquisición o uso de contingencia. & Usar parte de la reserva para una herramienta necesaria. \\
Calidad & Modifica criterios de aceptación, pruebas, estándares o nivel esperado del producto. & Agregar una revisión de navegación antes del cierre. \\
Riesgo & Incrementa o reduce la probabilidad o impacto de riesgos registrados. & Reducir alcance para bajar el riesgo de atraso. \\
Documentación & Requiere ajustar planes, matrices, registros o evidencias finales. & Cambiar presupuesto oficial de documentos relacionados. \\
\bottomrule
\end{tabularx}
}
\end{table}

\section{Roles y responsabilidades}

El control de cambios involucra a los integrantes internos y, cuando corresponde, al patrocinador o interesados externos. La participación depende del impacto del cambio solicitado.

\begin{table}[H]
\centering
\caption{Responsables del proceso de control de cambios}
\label{tab:responsables-cambios}
{\small
\begin{tabularx}{\textwidth}{p{4cm} Y}
\toprule
\textbf{Rol} & \textbf{Responsabilidad} \\
\midrule
Tamara Robles Camacho & Recibir solicitudes, coordinar revisión, validar impacto general y aprobar cambios menores o elevar cambios mayores. \\
Isaac Jiménez Blanco & Documentar solicitudes, revisar impacto en alcance, comunicaciones, calidad, riesgos y entregables relacionados. \\
Jorge Abarca Quiros & Apoyar el análisis de costos, riesgos, cronograma y consistencia documental. \\
Mauricio González Prendas & Evaluar impacto técnico en portal estudiantil, portal docente, portal administrativo, navegación y QA. \\
Carlos Sainz Arce & Evaluar impacto técnico en portal financiero, dashboard ejecutivo del sistema, navegación y QA. \\
Dr. Roberto Mora Fallas & Aprobar cambios relevantes que afecten alcance principal, costo total, fecha final o entregables estratégicos. \\
\bottomrule
\end{tabularx}
}
\end{table}

\section{Flujo del proceso de control de cambios}

El proceso se organiza en siete pasos secuenciales. Estos pasos buscan que cada cambio sea analizado antes de modificar el prototipo, el cronograma o la documentación.

\begin{table}[H]
\centering
\caption{Flujo del proceso de control de cambios}
\label{tab:flujo-cambios}
{\small
\begin{tabularx}{\textwidth}{p{1.4cm} p{3.3cm} Y}
\toprule
\textbf{Paso} & \textbf{Actividad} & \textbf{Descripción} \\
\midrule
1 & Identificación & Un integrante o interesado detecta una necesidad de modificación, problema o mejora. \\
2 & Registro de solicitud & Se completa una solicitud de cambio con fecha, solicitante, descripción, justificación y prioridad. \\
3 & Revisión preliminar & La directora del proyecto verifica si el cambio es válido, duplicado, urgente o fuera del alcance. \\
4 & Análisis de impacto & El equipo analiza efecto en alcance, cronograma, costo, calidad, riesgos, recursos y entregables. \\
5 & Decisión & El cambio se aprueba, rechaza, pospone o se aprueba con condiciones. \\
6 & Implementación & Si se aprueba, se actualiza el prototipo, cronograma, documentos o evidencia correspondiente. \\
7 & Comunicación y cierre & La decisión y el resultado se comunican al equipo y se deja trazabilidad final. \\
\bottomrule
\end{tabularx}
}
\end{table}

\section{Criterios de evaluación de impacto}

Cada solicitud debe evaluarse con los mismos criterios para evitar decisiones subjetivas. El análisis puede ser breve, pero debe cubrir las áreas principales del proyecto.

\begin{table}[H]
\centering
\caption{Criterios de evaluación de cambios}
\label{tab:criterios-evaluacion-cambios}
{\small
\begin{tabularx}{\textwidth}{p{3.5cm} Y}
\toprule
\textbf{Área} & \textbf{Preguntas de evaluación} \\
\midrule
Alcance & ¿Agrega o elimina módulos, pantallas, entregables o requisitos? ¿Contradice el alcance excluido? \\
Cronograma & ¿Afecta tareas, duraciones, dependencias, hitos o ruta crítica del cronograma en Microsoft Project? \\
Costo & ¿Usa presupuesto adicional, cambia rubros de adquisición o requiere aplicar contingencia? \\
Calidad & ¿Afecta criterios de aceptación, pruebas, navegación, usabilidad o consistencia documental? \\
Riesgos & ¿Aumenta o reduce riesgos como alcance creciente, retrasos, errores de navegación o baja evidencia final? \\
Recursos & ¿Sobrecarga a Tamara, Isaac, Jorge, Mauricio o Carlos? ¿Requiere apoyo externo? \\
Comunicación & ¿Debe informarse al patrocinador, stakeholders externos o todo el equipo? \\
\bottomrule
\end{tabularx}
}
\end{table}

\section{Decisiones posibles}

La evaluación de una solicitud puede producir cuatro decisiones. Esta clasificación permite mantener claridad y evitar cambios ambiguos.

\begin{table}[H]
\centering
\caption{Decisiones posibles ante una solicitud de cambio}
\label{tab:decisiones-cambios}
{\small
\begin{tabularx}{\textwidth}{p{3cm} Y}
\toprule
\textbf{Decisión} & \textbf{Descripción} \\
\midrule
Aprobado & El cambio se acepta y se implementa en el prototipo, documento o cronograma correspondiente. \\
Aprobado con condiciones & El cambio se acepta, pero con límites de alcance, fecha, costo o responsable. \\
Rechazado & El cambio no se implementa porque afecta demasiado el alcance, costo, tiempo o calidad. \\
Pospuesto & El cambio se considera valioso, pero se deja para una etapa posterior fuera del proyecto actual. \\
\bottomrule
\end{tabularx}
}
\end{table}

\section{Documentación y comunicación}

Toda solicitud aprobada debe dejar evidencia clara. Para cambios de bajo impacto puede bastar una anotación en el repositorio, una minuta o un mensaje formal de seguimiento. Para cambios que afecten alcance, cronograma, costo o calidad se debe utilizar la plantilla de solicitud de cambio y actualizar los documentos relacionados.

Los canales autorizados para comunicar cambios son GitHub, correo electrónico, documentos PDF, minutas de reunión y mensajes de coordinación cuando estos se usen solo como aviso. La mensajería instantánea no será suficiente para aprobar cambios importantes, aunque sí puede utilizarse para informar rápidamente que una solicitud debe ser revisada.

\section{Relación con otros documentos}

El proceso de control de cambios se relaciona con la definición del alcance, cronograma, plan de calidad, plan de comunicaciones, registro de riesgos, matriz de probabilidad e impacto, solicitudes de cambio y evaluación de cambios. Si un cambio modifica alguno de estos elementos, el documento correspondiente debe actualizarse para mantener la coherencia general del proyecto.

\section{Conclusión}

El Proceso de Control de Cambios permite proteger la estabilidad de la Plataforma Universitaria Integrada durante su desarrollo. Al establecer pasos claros para solicitar, analizar, aprobar, implementar y comunicar cambios, el equipo reduce el riesgo de alcance creciente, contradicciones documentales, atrasos y pérdida de calidad.

El proceso propuesto es suficientemente formal para mantener trazabilidad, pero también se adapta a la realidad académica del proyecto. De esta forma, el equipo puede reaccionar con agilidad ante necesidades relevantes sin perder control sobre el presupuesto de ₡17,625,200, el plazo planificado de tres meses, el alcance Front-End navegable y los criterios de aceptación definidos.

\end{document}
